%!TEX root = ../template.tex
%%%%%%%%%%%%%%%%%%%%%%%%%%%%%%%%%%%%%%%%%%%%%%%%%%%%%%%%%%%%%%%%%%%%
%% provisioning_esim.tex
%% NOVA thesis document file
%%
%% Chapter about provisioning the esim and connecting to the network
%%%%%%%%%%%%%%%%%%%%%%%%%%%%%%%%%%%%%%%%%%%%%%%%%%%%%%%%%%%%%%%%%%%%

\typeout{NT FILE provisioning_esim.tex}%

\chapter{Provisioning the eSIM and connecting to the network}
\label{ch:provisioning_esim}
In this chapter, we will provision an eSIM profile for the testbed and connect to the network. We will create an eSIM profile through a provisioning server and follow the steps to activate it on the \gls{UE}. Once the eSIM is activated, it can be used to connect to the network and transmit/receive data.

\section{Creating an eSIM profile}
\label{sec:creating_esim_profile}
To create an eSIM profile for the testbench, we will use Simlessly~\cite{simlessly}, specifically their \gls{RSP} tool, which allows custom eSIM profile creation. This platform requires an account to access the provisioning services.
Once logged in, the user is greeted with a dashboard, shown in Figure~\ref{fig:simlessly_dashboard}, where we can manage our eSIM profiles and view their status.

\begin{figure}[htbp]
    \centering
    \includegraphics[width=\textwidth]{Simlessly/simlessly_dashboard.png}
    \caption{Simlessly Dashboard.}
    \label{fig:simlessly_dashboard}
\end{figure}

From the dashboard, we can click the ``Generate AC'' button to create a new eSIM profile. In the prompt (Figure~\ref{fig:simlessly_generate_ac}), we fill in the details for the eSIM profile, provided in Table~\ref{tab:esim_profile_details}. 

\begin{figure}[htbp]
    \centering
    \includegraphics[width=0.6\textwidth]{Simlessly/GenerateAC.png}
    \caption{Generating an eSIM profile.}
    \label{fig:simlessly_generate_ac}
\end{figure}

\begin{table}[htbp]

    \centering
    \begin{tabular}{l|l}
        \hline
        \textbf{Key} & \textbf{Value} \\
        \hline
        ICCID &  8900101000000000001 \\
        \gls{IMSI} &  001010000000001 \\
        KI &  465B5CE8B199B49FAA5F0A2EE238A6BC \\
        OPC &  E8ED289DEBA952E4B4C389C94229606A \\
        \hline
    \end{tabular}
    \caption{eSIM profile details.}
    \label{tab:esim_profile_details}
\end{table}

Click on ``Generate'' to create the eSIM profile. Once generated, back on the dashboard, click on the QR code icon on the ``AC QR code'' column to display the Activation Code.

\begin{figure}[htbp]
    \centering
    \includegraphics[width=0.6\textwidth]{Simlessly/AC.png}
    \caption{eSIM Activation Code.}
    \label{fig:esim_activation_code}
\end{figure}

\section{Installing the profile on the \gls{UE}}
\label{sec:installing_esim_profile}
With the Activation Code displayed, the profile is installed on the phone by scanning it. On the \gls{UE}, we open the SIM manager, choose to add an eSIM and scan the QR code shown in Figure~\ref{fig:esim_activation_code}. The phone then asks for confirmation, downloads the profile from the provisioning server, and lists it in the SIM manager once it is installed. Figure~\ref{fig:esim_install} shows the three steps as they appear on the \gls{UE}.

\begin{figure}[htbp]
    \centering
    \includegraphics[width=0.3\textwidth]{Simlessly/lab_eSim1.jpeg}
    \hfill
    \includegraphics[width=0.3\textwidth]{Simlessly/lab_eSim2.jpeg}
    \hfill
    \includegraphics[width=0.3\textwidth]{Simlessly/lab_eSim3.jpeg}
    \caption[Installing the eSIM profile on the UE]{Installing the eSIM profile on the \gls{UE}. Left: the phone asks to confirm the \texttt{LAB\_eSIM} profile. Centre: the profile being downloaded from the provisioning server. Right: the SIM manager with the profile installed and enabled, and mobile data set to use it.}
    \label{fig:esim_install}
\end{figure}

Once the profile is enabled and selected for mobile data, the \gls{UE} attaches to the testbed's network as described in Appendix~\ref{app:testbed-setup}.